\SourceLeaf{17}
und bekanntlich sind $T,U,V$ der Reihe nach proportional zu den Kosinus der Winkel, die eine Normale zur Oberfläche mit den zu den Koordinatenachsen $x,y,z$ parallelen Geraden bildet, also der Winkel $QX,QY,QZ$. Daraus ergibt sich
\[
\begin{aligned}
\cos QX&=\frac{\pm T}{\BareRoot{(TT+UU+VV)}}\\
\cos QY&=\frac{\pm U}{\BareRoot{(TT+UU+VV)}}\\
\cos QZ&=\frac{\pm V}{\BareRoot{(TT+UU+VV)}}
\end{aligned}
\]
Es bleibt aber unklar, ob die oberen oder die unteren Vorzeichen zu nehmen sind. Um dies zu entscheiden, wählen wir auf der Geraden $PQ$, die in $P$ normal zur Oberfläche und \textit{nach außen} gerichtet ist, einen Punkt $P'$, der $P$ unendlich nahe liegt. Der Abstand sei $PP'=\mathrm{d}w$. Die Koordinaten des Punkts $P'$ sind dann der Reihe nach
\[
\begin{aligned}
x+\mathrm{d}w.\cos QX&=x+\mathrm{d}x\\
y+\mathrm{d}w.\cos QY&=y+\mathrm{d}y\\
z+\mathrm{d}w.\cos QZ&=z+\mathrm{d}z
\end{aligned}
\]
und daher ist der Zuwachs des Werts der Funktion $W$ vom Punkt $P$ (wo er $=0$ ist) bis zum Punkt $P'$
\[
\begin{aligned}
&=\mathrm{d}w.(T\cos QX+U\cos QY+V\cos QZ)\\
&=\pm\mathrm{d}w.\BareRoot{(TT+UU+VV)}
\end{aligned}
\]
Daraus folgt, dass die oberen Vorzeichen gelten, wenn die Funktion $W$ beim Verlassen des festen Teils des Körpers einen positiven und daher beim Eintritt in ihn einen negativen Wert annimmt. Im umgekehrten Fall gelten die unteren Vorzeichen. Unsere Oberfläche trennt nämlich den festen Teil des Körpers vom übrigen leeren Raum und zugleich die Raumteile, in denen $W$ einen positiven Wert annimmt, von denen, in denen die Funktion $W$ einen negativen Wert annimmt. Im Allgemeinen ist der Wert der Funktion $W$ \textit{entweder} außerhalb des Körpers positiv und innerhalb negativ, und dann sind die oberen Vorzeichen zu nehmen, \textit{oder} die Funktion $W$ ist außerhalb des Körpers negativ und innerhalb positiv, und dann gelten die unteren Vorzeichen.

Die Kosinus der übrigen Winkel, die wir in unseren Formeln benötigen, ergeben sich noch leichter. Wir haben nämlich
\[
\begin{aligned}
a&=x+r\cos MX\\
b&=y+r\sin MY\\
c&=z+r\sin MZ
\end{aligned}
\]
